\documentclass[11pt]{amsart}
\usepackage[T1]{fontenc}
\usepackage{lmodern}
\usepackage{amsmath,amssymb,amsthm,mathtools,booktabs}
\usepackage{xcolor}
\usepackage[colorlinks=true,linkcolor=blue!55!black,citecolor=blue!55!black,urlcolor=blue!55!black]{hyperref}

\newtheorem{theorem}{Theorem}[section]
\newtheorem{lemma}[theorem]{Lemma}
\newtheorem{proposition}[theorem]{Proposition}
\newtheorem{corollary}[theorem]{Corollary}
\theoremstyle{definition}
\newtheorem{definition}[theorem]{Definition}
\newtheorem{remark}[theorem]{Remark}

\title[Geometric Lexicographic Max-Cut on $K_7$]{The Continuous $n=7$ Case of Geometrically Weighted Lexicographic Max-Cut}
\author{Zijian Zeng}
\address{Institute of Computer Science and Digital Innovation, UCSI University, Kuala Lumpur 56000, Malaysia}
\email{zijianzeng@foxmail.com}
\email{1002266693@ucsiuniversity.edu.my}
\date{29 August 2026}
\subjclass[2020]{05C70}
\keywords{maximum cut, geometric edge weights, lexicographic order, Sturm sequence, exact certificate}
\hypersetup{pdftitle={The Continuous n=7 Case of Geometrically Weighted Lexicographic Max-Cut},pdfauthor={Zijian Zeng}}

\begin{document}

\begin{abstract}
For the complete graph $K_7$, order edges lexicographically and give the
$i$th edge weight $r^{21-i}$, where $1<r<2$.  We prove the $n=7$ case of
Mari\'c's global-optimality conjecture: throughout the full continuous
parameter interval, a maximum cut is one of the three initial-block cuts
$C_k=\{1,\ldots,k\}\mid\{k+1,\ldots,7\}$, $k=1,2,3$.  We isolate the two
transition parameters as unique positive roots of explicit integer
polynomials and classify all equality cases.  The proof is a finite exact
certificate comprising 185 rational Sturm-chain root counts and four
threshold identities.  A standard-library verifier independently
reconstructs all 64 cut polynomials and all 189 comparisons.  The result is
restricted to $n=7$ and does not establish the conjecture for larger $n$.
\end{abstract}

\maketitle

\section{Introduction}

Mari\'c introduced a geometrically weighted family of Max-Cut instances on
$K_n$ in which earlier edges in lexicographic order receive larger powers of
a parameter $r$~\cite{Maric2026}.  She proved an exact phase diagram within
the family of initial-block (``isolated'') cuts and conjectured that these
cuts are globally optimal for every $n\geq 7$ and $1<r<2$
~\cite[Theorem~4.10 and Conjecture~5.1]{Maric2026}.  The reported global
evidence exhausts all cuts for $n\leq21$ only on a grid of mesh $0.001$;
additional tests through $n=100$ consider a targeted near-isolated class.
The paper explicitly distinguishes this numerical evidence from a continuum
proof~\cite[Sections~3 and 5.1]{Maric2026}.

This note closes the first case, $n=7$, over the entire open interval
$(1,2)$.  The proof is computer-assisted but exact.  It neither samples the
parameter nor trusts a floating-point optimizer: all cut weights are integer
polynomials, threshold roots are isolated by rational signs, and every other
possible crossing is excluded by Sturm's theorem over rational arithmetic.
The accompanying verifier uses only the Python standard library and derives
the decisive quantities independently of the SymPy discovery program.

The new contribution is therefore global rather than restricted-family
optimality: we prove that the isolated cut selected by Mari\'c's phase diagram
beats every one of the 64 cuts of $K_7$ throughout its full parameter region,
and we classify all ties at the two algebraic transition points.

A targeted search of arXiv, DOI metadata, bibliographic indexes, and the
source paper's public code repository located no later public exact proof of the
continuous $n=7$ case through 29 August 2026.  This is a bounded bibliographic
report, not a claim about unpublished or unindexed work.

\section{Definitions and reductions}

Let $V=[7]$ and $N=\binom72=21$.  In lexicographic order the edges are
\[
(1,2),(1,3),\ldots,(1,7),(2,3),\ldots,(6,7),
\]
and the index of $(i,j)$, $i<j$, is
\[
 \operatorname{idx}(i,j)=7(i-1)-\binom{i}{2}+(j-i).
\]
The weight of $(i,j)$ is $r^{21-\operatorname{idx}(i,j)}$.  For a cut
$S\mid(V\setminus S)$ set
\[
 W_S(r)=\sum_{\substack{i<j\\ |S\cap\{i,j\}|=1}}
 r^{21-\operatorname{idx}(i,j)}.
\]
Complementary subsets give the same cut.  We therefore choose the unique
representative containing vertex 1.  Its other six membership bits enumerate
all $2^6=64$ cuts once, including the zero cut represented by $S=V$.

For $k=1,2,3$, define
\[
 C_k=\{1,\ldots,k\}\mid\{k+1,\ldots,7\}.
\]
Direct edge enumeration gives
\begin{align*}
W_{C_1}&=r^{20}+r^{19}+r^{18}+r^{17}+r^{16}+r^{15},\\
W_{C_2}&=r^{19}+r^{18}+r^{17}+r^{16}+r^{15}
 +r^{14}+r^{13}+r^{12}+r^{11}+r^{10},\\
W_{C_3}&=r^{18}+r^{17}+r^{16}+r^{15}+r^{13}+r^{12}
 +r^{11}+r^{10}+r^9+r^8+r^7+r^6.
\end{align*}
Put
\[
 q_2(x)=x^{13}+x^8-x^3-x^2-x-1,
 \qquad
 q_1(x)=x^{10}-x^4-x^3-x^2-x-1.
\]
Then
\begin{equation}\label{eq:threshold-identities}
 W_{C_2}-W_{C_3}=x^6q_2(x),\qquad
 W_{C_1}-W_{C_2}=x^{10}q_1(x).
\end{equation}

\begin{lemma}[Exact threshold isolation]\label{lem:thresholds}
The polynomial $q_2$ has a unique positive root $\alpha$, the polynomial
$q_1$ has a unique positive root $\beta$, and
\[
 1<\frac{539}{500}<\alpha<\frac{1079}{1000}
 <\frac{1229}{1000}<\beta<\frac{123}{100}<2.
\]
Both polynomials are negative before their positive root and positive after
it.
\end{lemma}

\begin{proof}
After zero coefficients are omitted, each polynomial has exactly one sign
change.  Descartes' rule gives at most one positive root.  Exact Horner
evaluation gives
\begin{align*}
500^{13}q_2(539/500)
 &=-1739575093438045166309424945633781,\\
1000^{13}q_2(1079/1000)
 &=24875693887459384063874283190637866839,\\
1000^{10}q_1(1229/1000)
 &=-15462253037576785810943471799,\\
100^{10}q_1(123/100)
 &=3331268605189126649.
\end{align*}
The intermediate value theorem proves existence in each stated interval;
Descartes' rule proves uniqueness.  Since each constant term is negative and
each leading coefficient is positive, the final sign assertion follows.
\end{proof}

\section{Main result}

\begin{theorem}[Continuous $n=7$ global optimality]\label{thm:main}
For every real $r\in(1,2)$ and every cut $S\mid(V\setminus S)$,
\[
W_S(r)\le
\begin{cases}
W_{C_3}(r),&1<r\le\alpha,\\
W_{C_2}(r),&\alpha\le r\le\beta,\\
W_{C_1}(r),&\beta\le r<2.
\end{cases}
\]
Up to complementation, $C_3$ is the unique maximizer on $(1,\alpha)$,
$C_2$ and $C_3$ are the only maximizers at $\alpha$, $C_2$ is unique on
$(\alpha,\beta)$, $C_1$ and $C_2$ are the only maximizers at $\beta$, and
$C_1$ is unique on $(\beta,2)$.
\end{theorem}

\section{Certificate or computational verification}

We record the exact finite proposition on which the theorem rests.  For a
target $C_k$ and a canonical competitor mask $m$, let
\[
 D_{k,m}(x)=W_{C_k}(x)-W_m(x).
\]

\begin{proposition}[Finite exact certificate]\label{prop:certificate}
For every nonadjacent competitor, the indicated difference is strictly
positive on the rational superinterval in the following table.  In the first
row, any exact $e$th-power factor of $x-1$ is removed before applying Sturm's
theorem.
\begin{center}
\begin{tabular}{ccccc}
\toprule
target & superinterval & Sturm cases & adjacent exceptions & $\#\{e>0\}$\\
\midrule
$C_3$ &$[1,1079/1000]$&62&$C_2$&34\\
$C_2$ &$[539/500,123/100]$&61&$C_3,C_1$&0\\
$C_1$ &$[1229/1000,2]$&62&$C_2$&0\\
\bottomrule
\end{tabular}
\end{center}
The 185 residual polynomials have degree at most 20, are positive at both
rational endpoints, and have no real root in their listed superinterval.
The only exceptional comparisons are the four directed forms of
~\eqref{eq:threshold-identities}.
\end{proposition}

\begin{proof}[Exact verification of Proposition~\ref{prop:certificate}]
All 189 directed difference polynomials are serialized coefficient by
coefficient in the certificate.
For each of the 185 nonexceptional residuals $f$, it constructs the rational
Sturm chain
\[
 f_0=f,\qquad f_1=f',\qquad
 f_{j+1}=-\operatorname{rem}(f_{j-1},f_j).
\]
If $V(t)$ is the number of sign changes at $t$ after zero entries are
discarded, direct rational evaluation gives $V(a)-V(b)=0$ at the listed
endpoints $a<b$.  Sturm's theorem gives no root in $(a,b)$; the positive
endpoint signs therefore give positivity throughout $[a,b]$.  Every removed
power of $x-1$ is positive on the theorem's domain $x>1$.

The independent verifier reconstructs the 21 edge exponents, all 64 weight
polynomials by two different evaluators, all 189 differences, every exact
division, and every Sturm chain using rational arithmetic.  It rejects
unknown fields, missing masks, altered coefficients, endpoint roots, and
nonzero root counts.  A clean run reports 185 Sturm cases, four threshold
identities, maximum residual degree 20, and maximum Sturm-chain length 21.
Thus the computation is an exact finite proof, not numerical evidence.
\end{proof}

\begin{proof}[Proof of Theorem~\ref{thm:main}]
Lemma~\ref{lem:thresholds} places the portion $1<x\leq\alpha$ inside the
first certified superinterval, $[\alpha,\beta]$ inside the second, and the
portion $\beta\leq x<2$ inside the third.
Proposition~\ref{prop:certificate} makes every nonadjacent
difference strictly positive there.

For the adjacent cuts,~\eqref{eq:threshold-identities} and
Lemma~\ref{lem:thresholds} show that
$W_{C_3}-W_{C_2}=-x^6q_2$ is positive before $\alpha$ and zero only there;
$W_{C_2}-W_{C_1}=-x^{10}q_1$ is positive before $\beta$ and zero only there.
Their opposite directed forms handle the phases to the right.  Hence the
claimed cut is globally optimal in each phase.  Strict positivity for every
nonexceptional competitor and uniqueness of the two positive threshold roots
give the equality classification.  Complementation only changes the chosen
representative, not the cut or its weight.
\end{proof}

\section{Equality cases and limitations}

The equality cases in Theorem~\ref{thm:main} are exhaustive among unordered
cuts.  The excluded endpoint $r=1$ has additional balanced-cut equalities and
is not part of the conjecture.  The certificate extends to $r=2$ for the
nonboundary comparisons, but the theorem deliberately retains the source
problem's open parameter range.

This result proves only the finite case $n=7$.  It does not establish a
factorization pattern or an induction capable of treating $n\geq8$; the
number of cuts grows exponentially, so the present certificate should not be
misrepresented as a proof of the mother conjecture.

\section*{Data, code, and computational disclosure}

The exact certificate is the file
\begin{center}
{\small\path{certificates/n7_sturm_certificate.json}.}
\end{center}
From the project root,
the decisive check is
\begin{verbatim}
python3 scripts/verify_certificate.py
\end{verbatim}
The discovery generator uses SymPy 1.14.0; the independent verifier uses only
Python's standard library.  No proof assistant was used.

All decisive mathematical claims are encoded in exact integer or rational data
and independently recomputed by the released verifier; no floating-point
output is used as a theorem.

\bibliographystyle{amsplain}
\bibliography{references}

\end{document}
